\documentclass[11pt, a4paper]{article}

\usepackage[T1]{fontenc}
\usepackage[utf8]{inputenc}
\usepackage{tgpagella}
\usepackage[spanish, es-noshorthands]{babel}
\usepackage{amsmath, amssymb, bm}
\usepackage{tabularx}
\usepackage{booktabs}
\usepackage[most]{tcolorbox}
\usepackage[margin=2.5cm]{geometry}
\usepackage{ragged2e}
\usepackage{fancyhdr}

\pagestyle{fancy}
\fancyhf{}
\setlength{\headheight}{16pt}
\lhead{\textbf{UCB Sede Tarija} | Ing. Mecatrónica}
\chead{}
\rhead{IMT-121 Circuitos Electrónicos I | BOM Lab RC}
\lfoot{Prof: M.Sc. Ing. Bernardo Quiroga Turdera}
\rfoot{Página \thepage}
\renewcommand{\headrulewidth}{0.4pt}

\definecolor{ucbblue}{RGB}{0, 51, 102}
\definecolor{warnred}{RGB}{200, 0, 0}

\begin{document}

\begin{tcolorbox}[colback=warnred!5!white, colframe=warnred, title=\textbf{BOM PROVISIONAL: Laboratorio de Transitorios RC}]
    \centering \textbf{Sujeto a Confirmación de Inventario y Equipo Institucional}
\end{tcolorbox}

\section*{1. Esquema de Rotación Propuesto}
Dado el inventario preliminar de \textbf{5 Osciloscopios} para aproximadamente \textbf{9 Grupos}:
\begin{itemize}
    \item \textbf{Rotación I:} Grupos 1 al 5 en estaciones de Osciloscopio. Grupos 6 al 9 en ensamblaje, medición manual con multímetro y cálculos teóricos.
    \item \textbf{Rotación II:} Intercambio total de puestos.
\end{itemize}

\section*{2. Lista de Materiales (BOM) Candidata}

\renewcommand{\arraystretch}{1.5}
\noindent
\begin{tabularx}{\textwidth}{>{\bfseries\RaggedRight}p{3cm} >{\centering\arraybackslash}p{2cm} >{\centering\arraybackslash}p{2cm} >{\RaggedRight\arraybackslash}X >{\centering\arraybackslash}p{2cm}}
\toprule
\textbf{Ítem} & \textbf{Cant. Grupo} & \textbf{Total (9 G)} & \textbf{Especificación Pendiente} & \textbf{Stock} \\ 
\midrule
Batería 9V & 1 & 9 & Tipo, conector y estado de carga real. & $\square$ Pendiente \\ \addlinespace
Conmutador SPDT & 1 & 9 & Verificación de contacto para evitar corto. & $\square$ Pendiente \\ \addlinespace
Resistor $10\,\text{k}\Omega$ & 1 & 9+ & Tolerancia, $1/4\,\text{W}$. & $\square$ Pendiente \\ \addlinespace
Resistor $22\,\text{k}\Omega$ & 1 & 9+ & Tolerancia, $1/4\,\text{W}$. & $\square$ Pendiente \\ \addlinespace
Capacitor Electrol. & 1 & 9+ & $470\,\mu\text{F}$. Voltaje máx $> 16\,\text{V}$. & $\square$ Pendiente \\ \addlinespace
Osciloscopio & - & \textbf{5 máx} & Confirmar sondas funcionales y disparo. & $\square$ Pendiente \\ \addlinespace
Multímetro Digital & 1 & 9 & Cables en buen estado. & $\square$ Pendiente \\ 
\bottomrule
\end{tabularx}

\vspace{0.8cm}
\begin{tcolorbox}[colback=gray!5, colframe=black, title=\textbf{Protocolo de Validación de Viabilidad}]
La sesión física no se confirmará hasta auditar:
\begin{enumerate}
    \item Tasa de Muestreo: Que el osciloscopio disponible posea bases de tiempo lo suficientemente amplias (segundos/división) para observar una carga lenta de $\sim 25$ a $\sim 50$ segundos.
    \item Aislamiento: Que la referencia de masa de las sondas no cortocircuite la alimentación del circuito RC.
\end{enumerate}
\end{tcolorbox}

\end{document}
